% Capitulo: Na prática: projetos de visualização
\chapter{Na prática: projetos de visualização}
\label{cap:proj-visualizacao}

Projetos de visualização têm uma armadilha própria: o produto final é tão visível que todo
mundo tem opinião sobre ele, e tão fácil de começar que a tentação é abrir a ferramenta e
sair montando gráficos. Este capítulo mostra como o modelo de Scrum com Kanban funciona num
projeto desse tipo, do pedido vago ao painel publicado, passando pelo que mais costuma dar
errado: construir antes de entender e entregar tudo de uma vez.

\section{Do pedido à pergunta}
\label{sec:pedido-pergunta}

Quase todo projeto de visualização começa com um pedido como \enquote{quero um dashboard de
vendas}. O pedido é legítimo, mas não é uma especificação. Ele não diz quem vai usar, com que
frequência, para decidir o quê, e o que conta como sucesso. Construir a partir dele é
apostar que o resultado vai coincidir com o que a pessoa tinha em mente, e a aposta quase
sempre perde.

O primeiro trabalho, portanto, é transformar o pedido em perguntas. Quem vai abrir o painel?
Que decisão essa pessoa toma, e com que frequência? O que ela faz hoje para tomar essa
decisão, e o que falta? Que número, se mudasse, faria ela agir de forma diferente? As
respostas costumam mostrar que o \enquote{dashboard de vendas} é, na verdade, a necessidade de
acompanhar três indicadores por região toda segunda-feira, com um alerta quando algum deles
sai da faixa normal. É um produto bem menor e bem mais útil do que o painel com trinta gráficos
que seria construído a partir do pedido original.

Essa conversa é também o momento de definir a frequência de atualização a partir do ritmo da
decisão (\veja{sec:ritmo-decisao}), de identificar o dono de cada métrica
(\veja{sec:donos}) e de escrever as primeiras regras que vão virar testes (\veja{sec:tdd}).
No quadro, ela corresponde à coluna de descoberta: nada vai para construção antes que essas
perguntas estejam respondidas.

\section{Protótipo e validação}
\label{sec:prototipo}

Com as perguntas respondidas, o passo seguinte ainda não é construir o painel de verdade. É
fazer um protótipo, o mais barato possível, e mostrá-lo a quem vai usar. Um rascunho em papel,
uma imagem montada numa ferramenta de desenho ou uma versão rápida na própria ferramenta de
BI, com dados fictícios ou com uma amostra, bastam. O objetivo é descobrir o que está errado
quando mudar ainda custa pouco.

A validação do protótipo deve ser feita com a pessoa que vai usar o painel, e não só com quem
pediu, porque nem sempre são a mesma pessoa. E deve ser conduzida com perguntas concretas:
olhando para este painel, que decisão você tomaria hoje? O que você procuraria primeiro? O que
está faltando? Quase sempre aparecem surpresas: o indicador que parecia central é ignorado, e
um detalhe que ninguém tinha mencionado é o que a pessoa procura primeiro.

As referências de visualização do capítulo~\ref{cap:self-service} ajudam a fazer protótipos
melhores desde o começo: uma tela só, poucos elementos, comparações sempre presentes e a
mensagem principal no título (\veja{sec:visualizacao}). A IA também ajuda aqui, gerando
rapidamente variações de layout ou dados fictícios realistas para o protótipo. O que ela não
faz é a conversa com quem vai usar.

\section{Entregas pequenas}
\label{sec:fatiar-dashboard}

O último erro comum é entregar o painel inteiro de uma vez, depois de semanas de construção.
Mesmo com protótipo validado, um painel grande entregue no fim acumula riscos: a fonte que
tinha um problema que ninguém viu, a regra que foi entendida de outro jeito, a mudança de
prioridade do negócio no meio do caminho.

A alternativa é fatiar. A unidade de entrega não é o painel, e sim a pergunta: a primeira
entrega responde à pergunta mais importante, com dados reais, testados e atualizados, mesmo
que o painel tenha só um gráfico. As entregas seguintes acrescentam as outras perguntas, uma de
cada vez, em ordem de valor. Cada fatia passa por todas as colunas do quadro, inclusive a
validação com o negócio, e por todos os itens da definição de pronto
(\veja{sec:definicao-pronto}).

Fatiar tem efeitos que vão além do risco. O usuário começa a usar o painel cedo, e o uso real
mostra o que importa mais do que qualquer conversa. Algumas fatias planejadas deixam de ser
necessárias, porque a primeira já resolveu o problema. E as métricas de uso, desde a primeira
entrega, indicam se o painel está sendo aberto na frequência imaginada
(\veja{sec:metricas-uso}).

\begin{dica}
Se um painel não pode ser dividido em entregas menores que façam sentido sozinhas, é
provável que ele esteja tentando responder perguntas demais. Vale voltar à descoberta.
\end{dica}

\rascunho{
  Relato de um projeto real de visualização, do pedido à entrega, sem nomes de empresa ou
  cliente: o pedido original, o que a descoberta revelou, o protótipo, como foi fatiado e o
  que aconteceu depois da primeira entrega. É o relato central do capítulo.
}

\begin{leituras}
  \leitura{few2006}{Para desenhar painéis que respondem a perguntas de relance, base de bons
    protótipos.}
  \leitura{knaflic2015}{Para conduzir a conversa com quem vai usar e para apresentar as
    entregas de forma que a mensagem fique clara.}
  \leitura{cohn2004}{Para escrever cada fatia como uma história com valor próprio para o
    usuário.}
\end{leituras}
